% !TeX root = ../../../../main.tex
% sections/chapter3/subsections/assumptions/market_structure.tex

\subsection{市場構造と主体の役割}
\label{sec:chapter3_assumptions_market_structure}
本モデルにおける市場構造と経済主体の基本的な役割について以下の3つの概念的な仮定をおく。

\paragraph{仮定1：家計は生産者を兼ねる}
\label{sec:chapter3_assumptions_market_structure_assumption1}
家計は生産者を兼ねる。

\paragraph{仮定2：家計は価格およびマクロ経済変数の受容者である}
\label{sec:chapter3_assumptions_market_structure_assumption2}
本モデルにおいて各家計の経済全体に占める規模は微小であると仮定する。
そのため個別の家計の行動が他の経済主体やマクロ経済全体に与える影響は軽微となる。
したがって家計は、自身が生産するものも含めた各生産者の財の価格や債券価格を所与として扱う「価格受容者」として行動する。
さらに家計は、経済全体の総体的な需給や政策によって決定される「マクロ経済変数」の受容者としても行動する。
このマクロ経済変数には、各種物価指数、名目利子率、名目為替レートといった集計的な価格体系のほか、マクロの総需要から受動的に決定される自身の生産などが該当する。
家計はこれらをすべて所与としたうえで消費および債券を最適化する。

\paragraph{仮定3：生産者はすべての家計の最適化行動を知っている}
\label{sec:chapter3_assumptions_market_structure_assumption3}
生産者としての家計は他のすべての家計の消費および債券の最適化行動について知ったうえで自身が生産する財の価格を最適化する。